% Scenario: notes / worked-gated   (tier: core)
% Exercises every path that DISCARDS worked content in a default (non-instructor)
% build, because each one typesets the hidden body into a throwaway \vbox that
% tagpdf will still walk unless tagging is suspended around it -- the phantom
% then nests inside the live surrounding paragraph as a <text-unit> within a
% <text>, i.e. PDF/UA Table 5 "P-Part" + "P-P".
%
% Three distinct layers gate this content, and each needed its own fix:
%   * didactic's \texlib@wrk@lead      -- \pf / \solution inside challenge+exercise
%   * didactic's own {solution} env    -- the standalone boxed solution
%   * thmenv's \texlib@worked@discard  -- the shared inline lead-in machinery
% Math inside the hidden bodies is deliberate: inline math is what turns into the
% empty <Formula> nodes that made the violation visible in the first place.
%
% The reference build is the STUDENT-facing one (no \ShowSolutions), which is
% exactly the copy that ships to a screen reader.
\documentclass{didactic}

\metasetup{
	institution    = {Smoke Test University},
	instructor     = {Test Instructor},
	season         = Fall,
	year           = 2026,
	course-subject = Math,
	course-number  = 101,
	course-title   = {Smoke Test Course},
	course-section = 1,
}

\meta{ unit = Lecture, number = 1, title = Gated Worked Content }

\begin{document}
\maketitle

\section{Gated Lead-Ins}
Each block below hides its worked content in this build.

\begin{exercise}
	Show that $x^3 - x - 1 = 0$ has a root in $(1, 2)$.
	\solution Let $f(x) = x^3 - x - 1$. Then $f(1) = -1 < 0$ and $f(2) = 5 > 0$,
	so the Intermediate Value Theorem gives a root in $(1, 2)$.
\end{exercise}

\begin{challenge}
	Generalize: for which $n$ does $x^n - x - 1$ have a real root?
	\pf For odd $n$, $x^n - x - 1 \to \pm\infty$ as $x \to \pm\infty$, so the
	Intermediate Value Theorem applies.
\end{challenge}

\section{Exposition Lead-Ins}
These stay visible: in example and question the lead-in is exposition, not a key.

\begin{example}
	Evaluate $\lim_{x \to 0} \frac{\sin x}{x}$.
	\solution The limit is $1$, by the squeeze theorem.
\end{example}

\begin{question}
	Is every continuous function differentiable?
	\ans No: $f(x) = |x|$ is continuous but not differentiable at $0$.
\end{question}

\section{Standalone Solution Box}
The boxed {solution} environment hides its body and leaves write-space.

\begin{exercise}
	Compute $\int_0^1 x^2 \, dx$.
\end{exercise}

% NO blank line before \begin{solution}: the paragraph above must still be OPEN
% when the discard box is created. That is the whole point -- with a blank line
% the environment starts its own block, the phantom lands as a sibling, and the
% document conforms even with the bug present. Opened mid-paragraph it nests, and
% an unguarded build fails Table 5 P-Part + P-P here. Do not "tidy" this gap in.
Evaluating it directly gives $\tfrac{1}{3}$, as the write-space below confirms.
\begin{solution}[2cm]
	$\int_0^1 x^2 \, dx = \left[ \tfrac{x^3}{3} \right]_0^1 = \tfrac{1}{3}$.
\end{solution}

\end{document}
